\originalchapter{Langevin、Fokker--Planck与临界动力学}
\label{ch:dynamics}
\source{8.334 L24--L25。}
\originalsection{Brownian粒子的平衡结构}
过阻尼粒子满足$\dot x=-\Gamma U'(x)+\eta(t)$，白噪声$\avg\eta=0$、$\avg{\eta(t)\eta(t')}=2D\delta(t-t')$。短时间增量$\Delta x=-\Gamma U'\Delta t+\sqrt{2D\Delta t}\,\zeta$，$\zeta$标准Gaussian；对任意平滑测试函数Taylor展开至$O(\Delta t)$，平均生成算符$\mathcal Lf=-\Gamma U'f'+Df''$。分部积分得到
\begin{equation}
\partial_tP=\partial_x[\Gamma U'P+D\partial_xP].
\end{equation}
概率流$J=-\Gamma U'P-D\partial_xP$。要求无流的$P_{\rm eq}\propto\e^{-\beta U}$，得到Einstein关系$D=\Gamma\kb T$。噪声与摩擦若不满足该关系，不能把稳态擅自称同温度的热平衡。

谐势$U=kx^2/2$的Ornstein--Uhlenbeck过程有$\avg{x(t)x(0)}=(\kb T/k)\e^{-\Gamma k|t|}$，响应$\chi(\omega)=\Gamma/(\Gamma k-\ii\omega)$，噪声谱$S_x=2D/[(\Gamma k)^2+\omega^2]$，直接满足经典涨落耗散$S_x=2\kb T\chi''/\omega$。这提供从随机方程到可测谱的完整例子。
\originalsection{非守恒序参量：Model A}
对场的物理自由能$F[\phi]$，取$\partial_t\phi=-\Gamma\delta F/\delta\phi+\zeta$，$\avg{\zeta(x,t)\zeta(x',t')}=2\Gamma\kb T\delta(x-x')\delta(t-t')$。泛函Fokker--Planck流同粒子情形，保证$P[\phi]\propto\e^{-\beta F}$。若Gaussian$F=\frac12\int[r\phi^2+K(\nabla\phi)^2]$，每$q$模衰减率$\Gamma(r+Kq^2)$，$\tau\sim r^{-1}\sim\xi^2$，Gaussian动态指数$z=2$。

响应$R(q,\omega)=[-\ii\omega+\Gamma(r+Kq^2)]^{-1}$，结构谱$C(q,\omega)=2\Gamma\kb T|R|^2$。频率积分给静态$\avg{|\phi_q|^2}=\kb T/(r+Kq^2)$，恢复前章Gaussian结果，这验证噪声归一及静动态一致性。
\originalsection{守恒序参量：Model B}
若$\int\phi$必须守恒，动力学应由电流$J=-\Gamma\nabla(\delta F/\delta\phi)+\eta_J$形成，$\partial_t\phi=\Gamma\nabla^2\delta F/\delta\phi-\nabla\cdot\eta_J$。噪声为电流噪声，谱额外$q^2$，不能继续用空间独立非守恒白噪声。

Gaussian衰减率$\Gamma q^2(r+Kq^2)$，临界$r=0$为$q^4$，$z=4$。相互作用情况下在标准Model B且无额外耦合的条件下守恒与静态异常维数联系为$z=4-\eta$；若序参量耦合动量、能量或其他慢量，可能属于其他动态普适类。相同静态Ising指数的两种模型可有不同弛豫时间。
\originalsection{动态RG与可观测时间}
动态尺度形式$C(q,\omega)=q^{-2+\eta-z}\mathcal C(\omega/q^z,q\xi)$，频率积分需给静态$q^{-2+\eta}$。$z$控制临界慢化，实验扫描速度若快于$\tau(\xi)$，就可能看到非平衡滞后而非平衡指数。局域Markov随机方程已忽略热库记忆，若噪声有长时间相关或惯性重要，需要扩展而不能只改$\Gamma$数值。

响应场路径积分可把噪声Gaussian概率积分，得到包含$\widetilde\phi(\partial_t\phi+\Gamma\delta F/\delta\phi)$与噪声二次项的动态作用量；其传播子是上述推迟响应，相关线为噪声乘两个响应。这个构造提供动态微扰图的起点，不能把静态传播子直接当频率依赖响应而忽略因果性。
